\documentclass[DIV=calc]{scrartcl}

\usepackage[utf8]{inputenc}
\usepackage[T1]{fontenc}
\usepackage[ngerman]{babel}
\usepackage{graphicx}
\usepackage[draft, markup=underlined]{changes}
\usepackage{csquotes}
\usepackage{ulem}
%\usepackage[dvipsnames]{xcolor}
\usepackage{paralist}
\usepackage{fixltx2e}
%\usepackage{ellipsis}
\usepackage[tracking=true]{microtype}
\usepackage{lmodern}                        % Ersatz fuer Computer Modern-Schriften
%\usepackage{hfoldsty}
%\usepackage{fourier}             % Schriftart
\usepackage[scaled=0.81]{helvet}     % Schriftart
\usepackage{url}
\usepackage{tocloft}           % Paket für Table of Contents
\usepackage{xcolor}
\definecolor{urlred}{HTML}{660000}
\usepackage{hyperref}
\hypersetup{
    colorlinks=true,    
    linkcolor=black,    % Farbe der internen Links (u.a. Table of Contents)
    urlcolor=black,    % Farbe der url-links
    citecolor=black} % Farbe der Literaturverzeichnis-Links
\usepackage{mdwlist} %Änderung Zeilenabstände itemize & enumerate
%\usepackage{draftwatermark} %Wasserzeichen ``Entwurf'' 
%\SetWatermarkText{Entwurf}

\parindent 0pt                 % Absatzeinrücken verhindern
\parskip 12pt                 % Absätze durch Lücke trennen
\setlength{\textheight}{23cm}
\usepackage{fancyhdr}
\pagestyle{fancy}
\fancyhead{} % clear all header fields
\cfoot{}
\lfoot{Zusammenkunft aller Physik-Fachschaften}
\rfoot{www.zapfev.de\\stapf@zapf.in}
\renewcommand{\headrulewidth}{0pt}
\renewcommand{\footrulewidth}{0.1pt}
\newcommand{\gen}{*innen}
\addto{\captionsngerman}{\renewcommand{\refname}{Quellen}}



%%%% Mit-TeXen Kommandoset
\usepackage[normalem]{ulem}
\usepackage{xcolor}
\newcommand{\replace}[2]{
    \sout{\textcolor{blue}{#1}}~\textcolor{blue}{#2}
}

\newcommand{\delete}[1]{
    \sout{#1}
}
\newcommand{\add}[1]{
    \textcolor{blue}{#1}
}
\newcommand{\okay}[1]{
    \textcolor{black}{#1}
}


\begin{document}
\hspace{0.87\textwidth}
\begin{minipage}{120pt}
\vspace{-1.8cm}
\includegraphics[width=80pt]{logo.pdf}
\centering
\small Zusammenkunft aller Physik-Fachschaften
\end{minipage}

\begin{center}
\huge{Resolution der Zusammenkunft aller Physik-Fachschaften}\vspace{.25\baselineskip}\\
\normalsize
\end{center}
\vspace{1cm}


\section*{Zur Unterstützung des Positionspapiers \glqq Damit Bildung keine Frage des Geldbeutels ist\grqq\ des BAföG-Bündnisses}

%\newpage


Die Zusammenkunft aller deutschsprachigen Physik-Fachschaften (ZaPF) begrüßt und unterstützt das Positionspapier \glqq Damit Bildung keine Frage des Geldbeutels ist\grqq\ des BAföG-Bündnisses in vollem Umfang. Um diese Unterstützung deutlich zu machen, wünscht die ZaPF, das genannte Positionspapier zu unterzeichnen.
    
\vspace*{\fill}
\begin{flushright}
Verabschiedet am 11.06.2019 in Bonn
\end{flushright}

\newpage
\begin{center}
\fbox{\includegraphics[height=0.9\textheight,page=1]{Positionspapier_buendnis.pdf}}
\fbox{\includegraphics[height=0.9\textheight,page=2]{Positionspapier_buendnis.pdf}}
\fbox{\includegraphics[height=0.9\textheight,page=3]{Positionspapier_buendnis.pdf}}
\end{center}

\end{document}